% Generated draft: requires mathematical and source review.
\documentclass[11pt]{article}
\usepackage[margin=25mm]{geometry}
\usepackage{fontspec}
\setmainfont{Latin Modern Roman}
\usepackage{amsmath,amssymb,mathtools,unicode-math}
\setmathfont{Latin Modern Math}
\usepackage{enumitem,needspace,xurl,hyperref}
\setlength{\parindent}{0pt}
\setlength{\parskip}{5pt}
\setlength{\emergencystretch}{3em}
\newcommand{\N}{\mathbb{N}}
\newcommand{\Z}{\mathbb{Z}}
\newcommand{\Q}{\mathbb{Q}}
\newcommand{\R}{\mathbb{R}}
\newcommand{\C}{\mathbb{C}}
\newcommand{\F}{\mathbb{F}}
\newcommand{\A}{\mathbb{A}}
\newcommand{\PP}{\mathbb{P}}
\newcommand{\E}{\mathbb{E}}
\newcommand{\eps}{\varepsilon}
\newcommand{\dd}{\,\mathrm{d}}
\begin{document}
\section*{UnsolvedMath ID: 2307013}
\textbf{Source problem code:} catalogue entry
\subsection{Definitions and mathematical statement}
A Riemann surface spread over the sphere consists of a Riemann surface $\mathcal S$ with a holomorphic projection $\pi:\mathcal S\to\widehat{\mathbb C}$, allowing ramification. If a simply connected $\mathcal S$ is parametrized conformally by a domain $U$, the composite projection is a meromorphic function $f:U\to\widehat{\mathbb C}$. This is a branched-surface setting, not necessarily an ordinary unbranched topological covering of the sphere. Value distribution concerns, for example, the counting of preimages of values and their Nevanlinna growth and deficiencies, in the appropriate uniformizing domain.

\textbf{Research problem (F. Huckeman, as printed).} Understand how the geometric structure of such a simply connected covering surface constrains the value distribution of its generating meromorphic function, extending the constituent-piece approach to include halfsheets with infinitely many branch points on their boundaries. The source does not define halfsheet or specify a particular bound, invariant, or characterization to prove.

The stated model is the surface generated by
\[
f(z)=e^z-z.
\]
Its critical points in the parameter plane occur at $z=2\pi i n$, $n\in\mathbb Z$, and are simple zeros of $f'(z)=e^z-1$, hence have local mapping degree two. Their corresponding projected values are $1-2\pi i n$. The source describes a right-half-plane constituent with these branch points on its boundary; distinguishing source points from projected values avoids treating the same coordinates as both. The term second-order branch point is retained in the local-degree-two reading, not as a claim of ramification order two in the convention that subtracts one.

The supplied summary reports this program open as of the 2018 edition, with no definitive resolution located.
\subsection{Short English statement}
Extend geometric value-distribution methods to halfsheet pieces whose boundaries contain infinitely many branch points, using the surface of $e^z-z$ as the supplied model. No specific quantitative theorem is stated.
\subsection{Sources}
\begin{itemize}
\item Hayman, \emph{Research Problems in Function Theory} (2018), Problem 7.13, attributed in the supplied text to F. Huckeman, \url{https://arxiv.org/abs/1809.07200}; source-TeX extraction, accessed 2026-07-29. The 2018 update was retained for status triage.
\item The source invokes Nevanlinna theory generally and supplies the example $e^z-z$, without a separate halfsheet-definition or value-distribution theorem citation.
\item Local registry \texttt{data/unsolvedmath/problems.json}, permanent ID $2307013$, source code catalogue entry, \emph{Research Problems in Function Theory --- Problem 7.13}. The record labels this open and OPEN-TRIAGE, with NEEDS\_REVIEW source-status and rights qualifications.
\end{itemize}
\end{document}
